
\documentclass{article}
\usepackage{colortbl}
\usepackage{makecell}
\usepackage{multirow}
\usepackage{supertabular}

\begin{document}

\newcounter{utterance}

\centering \large Interaction Transcript for game `cladder', experiment `full\_v1.5\_default', episode 9560 with llama{-}3.1{-}8b{-}instruct.
\vspace{24pt}

{ \footnotesize  \setcounter{utterance}{1}
\setlength{\tabcolsep}{0pt}
\begin{supertabular}{c@{$\;$}|p{.15\linewidth}@{}p{.15\linewidth}p{.15\linewidth}p{.15\linewidth}p{.15\linewidth}p{.15\linewidth}}
    \# & \multicolumn{2}{c}{Player} && \multicolumn{2}{c}{Game Master} \\
    \hline

    \theutterance \stepcounter{utterance}  
    & & & \multicolumn{4}{p{0.6\linewidth}}{
        \cellcolor[rgb]{0.9,0.9,0.9}{
            \makecell[{{p{\linewidth}}}]{
                \texttt{\tiny{[P1$\langle$GM]}}
                \texttt{You are an expert in causal inference. The following question is not a typical commonsense query, but rather a meticulously designed question created by a professor specializing in causal inference, intended to assess the students' mastery of the course content.} \\
\texttt{Imagine a self{-}contained, hypothetical world with only the following conditions, and without any unmentioned factors or causal relationships: Tijv has a direct effect on xevo and gyzp. Xevo has a direct effect on gyzp. We know that tijv causes not xevo. tijv and xevo causes gyzp. We observed an individual is not tijv. Would an individual is gyzp if not xevo instead of xevo?} \\
\texttt{Start your answer with 'Yes' or 'No', followed by additional reasoning or evidence to support your explanation.} \\
            }
        }
    }
     \\ \\

    \theutterance \stepcounter{utterance}  
    & \multicolumn{4}{p{0.6\linewidth}}{
        \cellcolor[rgb]{0.561,0.737,0.561}{
            \makecell[{{p{\linewidth}}}]{
                \texttt{\tiny{[P1$\rangle$GM]}}
                \texttt{No.} \\
\\ 
\texttt{To answer this question, we need to carefully analyze the given causal relationships and apply the principles of causal inference.} \\
\\ 
\texttt{We know that:} \\
\\ 
\texttt{1. Tijv has a direct effect on xevo and gyzp.} \\
\texttt{2. Xevo has a direct effect on gyzp.} \\
\texttt{3. Tijv causes not xevo (i.e., tijv is a cause of the absence of xevo).} \\
\texttt{4. Tijv and xevo both cause gyzp.} \\
\\ 
\texttt{Given that the individual is not tijv, we can infer that the absence of tijv is a cause of the absence of xevo (from point 3). However, the absence of tijv does not directly affect gyzp (from point 1).} \\
\\ 
\texttt{Now, we are asked to consider the scenario where the individual is not xevo instead of xevo. In this case, we need to examine the causal relationship between xevo and gyzp. Since xevo causes gyzp (from point 2), the absence of xevo would indeed affect gyzp.} \\
\\ 
\texttt{However, the key point to note is that the absence of tijv is a cause of the absence of xevo (from point 3). This means that the absence of xevo is not a direct effect of the absence of tijv, but rather an indirect effect through the causal relationship between tijv and xevo.} \\
\\ 
\texttt{In other words, the absence of tijv is a common cause of both the absence of xevo and the absence of gyzp. This is known as a "common cause" or "confounding" relationship.} \\
\\ 
\texttt{To resolve this, we can use the "backdoor criterion" from causal inference, which states that a path from a cause to an effect is blocked if there is a common cause that affects both the cause and the effect.} \\
\\ 
\texttt{In this case, the common cause is the absence of tijv, which affects both the absence of xevo and the absence of gyzp. Therefore, the path from the absence of xevo to gyzp is blocked by the common cause (the absence of tijv).} \\
\\ 
\texttt{As a result, we cannot conclude that the individual would be gyzp if not xevo instead of xevo. The absence of xevo would not directly affect gyzp, and the common cause (the absence of tijv) blocks the path from the absence of xevo to gyzp.} \\
            }
        }
    }
    & & \\ \\

    \theutterance \stepcounter{utterance}  
    & & & \multicolumn{2}{p{0.3\linewidth}}{
        \cellcolor[rgb]{0.9,0.9,0.9}{
            \makecell[{{p{\linewidth}}}]{
                \texttt{\tiny{[GM$|$GM]}}
                \texttt{No} \\
            }
        }
    }
    & & \\ \\

    \theutterance \stepcounter{utterance}  
    & & & \multicolumn{2}{p{0.3\linewidth}}{
        \cellcolor[rgb]{0.9,0.9,0.9}{
            \makecell[{{p{\linewidth}}}]{
                \texttt{\tiny{[GM$|$GM]}}
                \texttt{no} \\
            }
        }
    }
    & & \\ \\

    \theutterance \stepcounter{utterance}  
    & & & \multicolumn{2}{p{0.3\linewidth}}{
        \cellcolor[rgb]{0.9,0.9,0.9}{
            \makecell[{{p{\linewidth}}}]{
                \texttt{\tiny{[GM$|$GM]}}
                \texttt{game\_result = WIN} \\
            }
        }
    }
    & & \\ \\

\end{supertabular}
}

\end{document}
